\textbf{Phần 1. Trắc nghiệm lựa chọn}
\setcounterpageref{ex}{0}

\begin{ex} %[2D12V63]
Giá trị của biểu thức $456 + 234 - 190$ là
\choice
{$490$}
{$510$}
{\True $500$}
{$520$}
\loigiai{
$456 + 234 = 690$; $690 - 190 = 500$.
\textbf{Đáp án của câu hỏi là C}
}
\end{ex}

\begin{ex} %[2D12V63]
Tổng của hai số là $750$. Nếu số hạng thứ nhất tăng $50$ và số hạng thứ hai giảm $20$ thì tổng mới là
\choice
{$720$}
{\True $780$}
{$800$}
{$770$}
\loigiai{
Tổng mới: $750 + 50 - 20 = 780$.
\textbf{Đáp án của câu hỏi là B}
}
\end{ex}

\begin{ex} %[2D12V63]
Tìm $x$ biết: $x + 145 = 500 - 120$.
\choice
{$245$}
{$225$}
{$380$}
{\True $235$}
\loigiai{
$500 - 120 = 380$. $x = 380 - 145 = 235$.
\textbf{Đáp án của câu hỏi là D}
}
\end{ex}

\begin{ex} %[2D12V63]
Một đoàn tàu chở $680$ hành khách. Đến ga thứ nhất có $125$ người xuống tàu và $45$ người lên tàu. Trên tàu còn lại bao nhiêu hành khách?
\choice
{\True $600$ người}
{$590$ người}
{$610$ người}
{$555$ người}
\loigiai{
$680 - 125 + 45 = 600$ (hành khách).
\textbf{Đáp án của câu hỏi là A}
}
\end{ex}

\begin{ex} %[2D12V63]
Điền dấu thích hợp: $450 + 350 \dots 1 000 - 200$.
\choice
{$>$}
{$<$}
{\True $=$}
{Không so sánh được}
\loigiai{
$450 + 350 = 800$; $1 000 - 200 = 800$. Hai vế bằng nhau.
\textbf{Đáp án của câu hỏi là C}
}
\end{ex}

\begin{ex} %[2D12V63]
Một người có hai tấm vải: tấm thứ nhất dài $245\text{ m}$, tấm thứ hai dài hơn tấm thứ nhất $30\text{ m}$. Cả hai tấm vải dài bao nhiêu mét?
\choice
{$510\text{ m}$}
{\True $520\text{ m}$}
{$530\text{ m}$}
{$275\text{ m}$}
\loigiai{
Tấm 2: $245 + 30 = 275$ m. Cả hai tấm: $245 + 275 = 520$ m.
\textbf{Đáp án của câu hỏi là B}
}
\end{ex}

\begin{ex} %[2D12V63]
Số tự nhiên lớn nhất có ba chữ số khác nhau trừ đi số tự nhiên bé nhất có ba chữ số khác nhau được kết quả là
\choice
{$888$}
{$899$}
{$875$}
{\True $885$}
\loigiai{
Số lớn nhất khác nhau: 987. Bé nhất khác nhau: 102. Hiệu: $987 - 102 = 885$.
\textbf{Đáp án của câu hỏi là D}
}
\end{ex}

\begin{ex} %[2D12V63]
Thùng thứ nhất có $350$ lít dầu, thùng thứ hai có ít hơn thùng thứ nhất $70$ lít. Cả hai thùng có bao nhiêu lít dầu?
\choice
{\True $630$ lít}
{$620$ lít}
{$640$ lít}
{$560$ lít}
\loigiai{
Thùng 2: $350 - 70 = 280$ l. Cả hai thùng: $350 + 280 = 630$ l.
\textbf{Đáp án của câu hỏi là A}
}
\end{ex}

\begin{ex} %[2D12V63]
Tính tổng của ba số tròn trăm liên tiếp, biết số bé nhất là $200$.
\choice
{$800$}
{$700$}
{\True $900$}
{$1 000$}
\loigiai{
Ba số là 200, 300, 400. Tổng: $200 + 300 + 400 = 900$.
\textbf{Đáp án của câu hỏi là C}
}
\end{ex}

\begin{ex} %[2D12V63]
Bác Nam nuôi $320$ con gà và $280$ con vịt. Sau khi bán đi một số con gia cầm thì bác còn lại $150$ con. Bác Nam đã bán bao nhiêu con gia cầm?
\choice
{$440$ con}
{\True $450$ con}
{$460$ con}
{$350$ con}
\loigiai{
Tổng gia cầm: $320 + 280 = 600$ con. Bác đã bán: $600 - 150 = 450$ con.
\textbf{Đáp án của câu hỏi là B}
}
\end{ex}

\begin{ex} %[2D12V63]
Một hình tam giác có chu vi $600\text{ cm}$. Cạnh thứ nhất dài $180\text{ cm}$, cạnh thứ hai dài $220\text{ cm}$. Độ dài cạnh thứ ba là
\choice
{$190\text{ cm}$}
{$210\text{ cm}$}
{$250\text{ cm}$}
{\True $200\text{ cm}$}
\loigiai{
Cạnh thứ ba: $600 - (180 + 220) = 200$ cm.
\textbf{Đáp án của câu hỏi là D}
}
\end{ex}

\begin{ex} %[2D12V63]
Cho $x + y = 500$. Nếu giữ nguyên $x$ và bớt $y$ đi $120$ đơn vị thì tổng mới là bao nhiêu?
\choice
{\True $380$}
{$420$}
{$620$}
{$480$}
\loigiai{
Tổng mới: $500 - 120 = 380$.
\textbf{Đáp án của câu hỏi là A}
}
\end{ex}

\textbf{Phần 2. Trắc nghiệm đúng sai}
\setcounterpageref{ex}{0}

\begin{ex} %[2D12V63]
Xét tính đúng sai của các bài toán tổng hợp cộng trừ phạm vi 1 000:
\choiceTF
{\True 350 + 450 - 200 = 600.}
{1 000 - 300 - 200 = 600.}
{\True Nếu a + b = c thì c - a = b.}
{\True Tổng của số chẵn lớn nhất có 3 chữ số và số lẻ bé nhất có 1 chữ số bằng 999.}
\loigiai{
\begin{itemize}
    \item Ý a) Đúng: 800 - 200 = 600. Đúng.
    \item Ý b) Sai: 1 000 - 500 = 500, không phải 600. Sai.
    \item Ý c) Đúng: Tính chất của phép trừ và phép cộng. Đúng.
    \item Ý d) Đúng: 998 + 1 = 999. Đúng.
\end{itemize}
\begin{itemize}
    \item a) \textbf{Đúng}.
    \item b) \textbf{Sai}.
    \item c) \textbf{Đúng}.
    \item d) \textbf{Đúng}.
\end{itemize}
}
\end{ex}

\begin{ex} %[2D12V63]
Cửa hàng vật liệu có 900 bao xi măng. Buổi sáng bán 350 bao, buổi chiều bán 250 bao:
\choiceTF
{\True Cả ngày bán được 600 bao xi măng.}
{\True Cửa hàng còn lại 300 bao xi măng.}
{Số bao bán buổi sáng gấp đôi buổi chiều.}
{\True Số xi măng còn lại đúng bằng 1/2 số xi măng đã bán.}
\loigiai{
\begin{itemize}
    \item Ý a) Đúng: 350 + 250 = 600 bao. Đúng.
    \item Ý b) Đúng: 900 - 600 = 300 bao. Đúng.
    \item Ý c) Sai: 250 x 2 = 500 != 350. Sai.
    \item Ý d) Đúng: 300 = 600 : 2. Đúng.
\end{itemize}
\begin{itemize}
    \item a) \textbf{Đúng}.
    \item b) \textbf{Đúng}.
    \item c) \textbf{Sai}.
    \item d) \textbf{Đúng}.
\end{itemize}
}
\end{ex}

\textbf{Phần 3. Trả lời ngắn}
\setcounterpageref{ex}{0}

\begin{ex} %[2D12V63]
Tính giá trị: $523 + 277 - 350$.
\shortans{450}
\loigiai{
\textbf{Đáp số:} $450$\\
$800 - 350 = 450$.
}
\end{ex}

\begin{ex} %[2D12V63]
Tìm $x$ biết: $x - 175 = 625 - 200$.
\shortans{600}
\loigiai{
\textbf{Đáp số:} $600$\\
$625 - 200 = 425$. $x = 425 + 175 = 600$.
}
\end{ex}

\begin{ex} %[2D12V63]
Kho gạo có 1 000 kg gạo. Sau hai lần xuất thì còn lại 320 kg. Hỏi kho đã xuất bao nhiêu ki-lô-gam gạo?
\shortans{680}
\loigiai{
\textbf{Đáp số:} $680$\\
$1 000 - 320 = 680$ kg.
}
\end{ex}

\begin{ex} %[2D12V63]
Tổng của hai số là 900. Số thứ nhất là 345. Tìm số thứ hai.
\shortans{555}
\loigiai{
\textbf{Đáp số:} $555$\\
$900 - 345 = 555$.
}
\end{ex}

\begin{ex} %[2D12V63]
Một hình tam giác có độ dài các cạnh lần lượt là 250 cm, 300 cm và 450 cm. Chu vi hình tam giác đó là bao nhiêu mét?
\shortans{10}
\loigiai{
\textbf{Đáp số:} $10$\\
Chu vi: $250 + 300 + 450 = 1 000\text{ cm} = 10$ m.
}
\end{ex}

\begin{ex} %[2D12V63]
Có 3 xe chở hàng: xe 1 chở 280 bao, xe 2 chở 320 bao, xe 3 chở 400 bao. Cả 3 xe chở bao nhiêu bao?
\shortans{1 000}
\loigiai{
\textbf{Đáp số:} $1 000$\\
$280 + 320 + 400 = 1 000$ bao.
}
\end{ex}
